% JK: AJE Prac of Epi abstracts: 200 words unstructured
Average total durations selling or buying sex
are key determinants of HIV/STI transmission dynamics and service coverage.
Such durations are usually estimated from
cross-sectional surveys of people actively selling/buying sex.
However, cross-sectional surveys tend to
over-sample individuals selling/buying sex for longer,
while survey-reported durations will be right-censored.
Prior works have not precisely addressed these two competing biases.
In this paper, we derive the exact relationship between
the distribution of durations reported in cross-sectional surveys and
the distribution of total durations in the population overall.
These distributions will be identical
if total durations are exponentially distributed;
otherwise, reported durations will typically underestimate the variance in total durations.
We then develop a parametric Bayesian model of this relationship,
and apply it to re-analyze aggregate data from a global systematic review.
Our model-estimated average total durations are
significantly shorter than previously published for all regions --- \eg
% JK: I chose these regions to highlight b/c they have the most data
in Africa: 3.7 (\CI: 3.4--3.9) versus 5.5 years;
in Asia: 1.4 (1.2--1.5) versus 2.8--4.0 years.
Shorter average durations then suggest that
HIV incidence in sex work could be higher than previously thought,
to match the same HIV prevalence data.
Our approach can be applied beyond sex work,
to characterize duration distributions for any transient risk/disease state.
